\documentclass[a4paper,11pt]{ltjsarticle}
\usepackage[margin=25mm]{geometry}
\usepackage{luatexja-fontspec}
\setmainjfont{Hiragino Mincho ProN}
\setsansjfont{Hiragino Sans}
\usepackage{listings,enumitem,tabularx,longtable}
\lstset{basicstyle=\small\ttfamily,numbers=left,numberstyle=\tiny,frame=single,breaklines=true,xleftmargin=2em,aboveskip=5pt,belowskip=4pt}
\begin{document}
\section{折板の寸法}
測定台の左端を起点とした。\\
同じ手順でもう一組の板を読む。
\begin{lstlisting}
ridge = 12
valley = 9
span = ridge + valley
\end{lstlisting}
\begin{description}[style=nextline,leftmargin=4em,labelwidth=3em]
\item[峰] 板の外側の折り線。
\item[試験用の長い谷線] 内側に引いた細い線。
\end{description}
\begin{enumerate}
\item 切片を並べる。
\item 頂点を探す。
\begin{itemize}
\item 白印の位置を測る。
\item 黒印の位置を測る。
\end{itemize}
\end{enumerate}
\noindent\begin{tabularx}{\linewidth}{p{32mm}p{57mm}X}
測定片 & 開口の向き & 残りの幅 \\
第一片 & 北側 & 17 \\
第二片 & 南側 & 23
\end{tabularx}
\begin{longtable}{lp{35mm}r}
番号 & 説明 & 回数 \\
一 & 端を押さえて静かに折り目を開く & 3 \\
二 & 折り目を閉じる & 5
\end{longtable}
\noindent\begin{minipage}[t]{.46\linewidth}\subsection*{手前の台}細い方の板を置く。\end{minipage}\hfill
\begin{minipage}[t]{.46\linewidth}\subsection*{奥の台}広い方の板を置く。\end{minipage}
\end{document}
